\section{Introduction}

The scattering problem for general nonlinear
f\/ield equations has been intensively studied for many years with
considerable progress, a seminal result being the establishment of
asymptotic completeness in energy space for the
nonlinear Schr{\"o}dinger equation in space dimension \mbox{$n\geq 3$} with
interaction $|\varphi|^{p-1}\varphi$, where
$1+4/n\leq p\leq 1+4/(n-2)$, and analogous results for the nonlinear Klein--Gordon
equation~\cite{GiVe}. Still, many important questions remain
unanswered, in particular relating to the existence of wave operators
def\/ined on appropriate subspaces of scattering states and to
asymptotic completeness for more general interactions, although many
partial results have been obtained. We refer to~\cite{Strauss,Ginibre} for an overview.
One source of dif\/f\/iculties encountered can be traced back to the
singular nature of pointwise
multiplication of functions with respect to appropriate Lebesgue or
Sobolev norms. It therefore appears natural to look for f\/ield equations where
this part of the problem can be eliminated while preserving as much
as possible of the remaining structure. One such possibility
is to deform the standard product appropriately, and  our
purpose  in the present article is to exploit this idea.

Deformations of the kind mentioned applied to relativistic f\/ield
equations have turned out to be of relevance for understanding
non-perturbative aspects of string theory, see e.g.~\cite{GMS,HKL,DJN}.
We shall limit ourselves to studying a deformed version of the
nonlinear Schr{\"o}dinger equation in an even number $n=2d$ of space
dimensions. The rather crude Hilbert space
techniques~\cite{Reed,RS} we apply allow us to consider polynomial
interactions only. We hope that an application of more modern techniques
\cite{Caz} will allow a treatment of a larger class of interactions.
Our main purpose will be to demonstrate that, for
the deformed equation,  the Cauchy problem is in a certain sense more
regular as compared to the
classical equation as well as to set up a natural scattering framework,
 including appropriate decay estimates of the free time evolution
and construction of wave operators under rather mild conditions on the
interaction.

The deformed, or
noncommutative, version of the nonlinear Schr{\"o}dinger
equation (NCNLS) in $2d$ space dimensions can be written as
\begin{gather}
\label{NLSch}
\big(i\partial_t-{\Delta}\big)\varphi(x,t) =
\varphi\,\star_\theta F_{\star_\theta}(\varphi^*\star_\theta\varphi)(x,t) ,
\end{gather}
where $\star_\theta$ denotes the Moyal product (see below) of functions of
$2d$ space variables $x$, $F_{\star_\theta}$ denotes a real polynomial
with respect to this product and ${\Delta}=-\sum_{i=1}^{2d}\partial_i^2$ is
the standard Laplacian in~$2d$ variables.
The Moyal product considered here is given by
\begin{gather*}
%\label{mop}
f\star_\theta g(x):=(2\pi)^{-2d}\int e^{-iyz} f\left(x-\frac\Theta2
y\right)  g(x+z)\,d^{2d} y\,d^{2d} z .
\end{gather*}
The constant skew-symmetric $(2d\times2d)$-matrix $\Theta$ is assumed to
be given in the canonical form
\[
\Theta  =  \theta \begin{pmatrix} 0 & I_d\\ -I_d & 0\end{pmatrix} ,
\]
where $I_d$ denotes the $d\times d$ identity matrix and
$\theta>0$ is called the deformation parameter.

Def\/ining
$\varphi_\theta(x,t) = \varphi(\theta^\frac 12 x,\theta t)$,
we have the scaling identity
\[
(\varphi\star_\theta \psi)_{\theta} = \varphi_{\theta}\star_1 \psi_{\theta}.
\]
It follows that $\varphi$ satisf\/ies (\ref{NLSch}) if and only if
\begin{gather*}
%\label{NLSch1}
\big(i\partial_t-\Delta\big)\varphi_\theta(x,t) =
\theta \varphi_\theta \star F_\star(\varphi_\theta^*\star\varphi_\theta)(x,t) ,
\end{gather*}
where we have dropped the subscript on the $\star$-product when $\theta =1$.


The $\star$-product is intimately connected to the so-called Weyl
quantization map $W$. This map associates an operator
$W(f)$ on $L^2(\mathbb R^d)$ to appropriate functions (or distributions) $f$ on $\mathbb R^{2d}$
such that the kernel $K_W(f)$ of $W(f)$ is given by
\begin{gather*}
%\label{W}
K_W(f)(x,y)=(2\pi)^{-d}\int_{\mathbb
R^d}f\left(\frac{x+y}{2},p\right)e^{i(x-y)\cdot p}\,dp  .
\end{gather*}
 It is easily seen that $W$ is an
isomorphism from $L^2(\mathbb R^{2d})$ onto the Hilbert space $\cal H$ of
Hilbert--Schmidt operators on $L^2(\mathbb R^d)$ fulf\/illing
\begin{gather}
\label{Wunitary}
\Vert W(f)\Vert_2 = (2\pi)^{-d/2}\Vert f\Vert_{L^2(\mathbb R^{2d})} ,
\end{gather}
where $\Vert\cdot\Vert_2$ denotes the Hilbert--Schmidt norm. Moreover,
$W$ maps the space ${\cal S}(\mathbb R^{2d})$ onto the space of
operators whose kernel is a Schwartz function and its relation to the
$\star$-product is exhibited by the identity
\[
W(f\star g) = W(f)W(g) .
\]
Further properties of the $\star$-product can be found in
e.g.~\cite{GGBISV}.

By use of $W$, equation~\eqref{NLSch} can be restated (see also \cite{DJN})
as a dif\/ferential equation
\begin{gather}
\label{NLSch2}
i\partial_t\phi - 2 \sum_{k=1}^d[a_k^*,[a_k,\phi]] =\theta \phi F(\phi^*\phi) ,
\end{gather}
 for  the operator-valued function $\phi(t) := W(\varphi_\theta(\cdot,t))$,
where we have introduced the creation and annihilation operators
\[
a_k =\frac{1}{\sqrt 2}(x_k + \partial_k)\qquad\mbox{and}\qquad a_k^* =
\frac{1}{\sqrt 2}(x_k - \partial_k) .
\]
The operator
\[
{\bf \Delta} = 2 \sum_{k=1}^d \mbox{ad}\, a_k^*\;\mbox{ad}\, a_k,
\]
with domain $D({\bf\Delta})$,
is def\/ined in a natural way as a self-adjoint operator on $\cal H$ by the
relations
\begin{gather}
\label{tria}
{\bf \Delta} =  W\Delta W^{-1} ,\qquad D({\bf \Delta})=WD(\Delta) ,
\end{gather}
where $\Delta$ denotes the standard self-adjoint $2d$-dimensional Laplace
operator with maximal domain $D(\Delta)=H_2^2(\mathbb R^{2d})$.

We shall primarily be interested in globally def\/ined \emph{mild}
solutions to the Cauchy problem associated to
equation (\ref{NLSch2}), that is continuous solutions $\phi:\mathbb
R\to{\cal H}$ to the corresponding integral equation
\begin{gather}
\label{NLSch3}
\phi(t) = e^{-i(t-t_0){\bf \Delta}}\phi_0 - i\int_{t_0}^t e^{i(s-t){\bf \Delta}}
\phi(s) F\big(|\phi(s)|^2\big)\,ds .
\end{gather}
This equation is weaker than (\ref{NLSch2}) in the sense that if
$\phi: I\to D({\bf \Delta})$ is a continuously dif\/ferentiable solution to
(\ref{NLSch2}) def\/ined on an interval $I$ containing $t_0$, i.e.\ a
\emph{strong solution} on $I$, then it also fulf\/ills (\ref{NLSch3}).
This latter equation  naturally f\/its into
the standard Hilbert space framework for evolution equations, see
e.g.~\cite{RS,Reed}. The following theorem is proven in Section~\ref{secth1}
below.

